\documentclass[11pt,a4paper]{article}
\usepackage[margin=2.2cm]{geometry}
\usepackage[T1]{fontenc}
\usepackage{amsmath,amssymb,xcolor,enumitem,booktabs,graphicx,tabularx}
\usepackage[normalem]{ulem}
\setlength{\parskip}{0.6em}\setlength{\parindent}{0pt}
\definecolor{newc}{HTML}{1f6f43}
\newcommand{\new}[1]{\textcolor{newc}{#1}}
\newcommand{\old}[1]{\textcolor{gray}{\sout{#1}}}
\newcommand{\note}[1]{{\small\itshape [#1]}}
\newcommand{\lab}[1]{\textsf{\footnotesize\detokenize{#1}}}
\renewcommand{\ref}[1]{\lab{#1}}
\renewcommand{\eqref}[1]{(\lab{#1})}
\newcommand{\citep}[1]{[\lab{#1}]}
\newcommand{\citet}[1]{\lab{#1}}
\newcommand{\mockup}[1]{\fbox{\small mockup: #1}}
\title{Proposed change, ``A reference point''}
\date{24 November 2026, version 2}
\begin{document}\sloppy
\maketitle

\section*{How to read this}
This answers what is left of \texttt{method\_reference\_comments.md} for
\lab{sec:method:reference} in \texttt{4.Implementation.tex}. The lock (B2) was
done in [A109]. Agreed on 24 November: the role told by design, not by
results (A4); the sentence on A$^\star$, ICBS and ICTS dropped (B4); no figure
(E1). The rest (A3, B3, C1 to C5, D1, D2, D4) is fixed as the comments say.

New in version 2 (agreed): the section is split. How Token Passing works in
the loop stays in the Method, under the plain name ``Token Passing''. Its role
in the experiment (what it is measured against, the values it sets, how its
gain is read) moves to the Experiments, ``Factors'', where a mockup held its
place. The label \lab{sec:method:reference} stays, so no reference breaks.
B1, why Token Passing and not its variant with task swaps, is a proposal (see
the question). No behaviour of Token Passing changes.

Black text stays, \new{green} is new and \old{struck} text goes.
\note{Notes in brackets} say why. Nothing below is in the thesis yet.

\section*{Part 1. What was checked}
\begin{itemize}[nosep]
  \item \texttt{planners/token\_passing.py}: agents ask in agent id order
    (``Ma et al.\ leave the order open''); the nearest available task by hop
    distance; a task is available if no other path ends at its pickup or
    delivery; with no task an agent stays, unless a waiting task delivers
    where it stands, then it goes to the nearest free endpoint that is no
    shelf or delivery point; each path is a space-time A$^\star$ path
    against the token; a vertex is taken one step before arrival until one
    step after leaving.
  \item Ma et al.\ (2017) give two solvers, Token Passing and Token Passing
    with Task Swaps, where an agent may take over a task from another agent
    that would reach it later. Van Knippenberg et al.\ do not use Token
    Passing. A$^\star$ is not in \texttt{references.bib}.
  \item Appendix~E: a load sweep of Token Passing under fixed storage sets
    the main load of each fleet, chosen ``because it needs no training''.
    Appendix~G: sweeps of Token Passing set the storage period and the
    pallet size. The Results run it under every storage rule at the main
    load.
  \item The Discussion opens with ``only for Token Passing can the gain be
    measured''. That is a result, so it stays there.
  \item The Problem Formulation already has the scope item ``One agent per
    workstation at a time'' (\lab{sec:pf:scope}).
\end{itemize}

\section*{Part 2. The Method opening}

\ldots Free learned agents are paired with waiting tasks by the assignment
rule, and Token Passing instead selects its own tasks. Each \new{learned}
agent then chooses one move from what it sees
(Table~\ref{tab:information})\new{, and a Token Passing agent takes the next
step of its planned path}. \ldots \note{A3, step 3 did not fit Token
Passing.}

\ldots They end with the reward, the training and \old{a reference point
outside the main comparison}\new{Token Passing}. \note{The section's new
name.}

\section*{Part 3. The Method section as it would read}

\subsection*{\old{A reference point}\new{Token Passing}}

\old{Token Passing is run next to the learned controllers but is not part of
the comparison the research questions make. It checks that the learned
controllers can compete with a classic planner. It is run at the main loads,
not across the full design. Token Passing runs under every storage rule, so
its gain can be set next to the learned controllers'.}
\note{A4. The role moves to the Experiments, Part 4.}

\new{Token Passing \citep{ma2017lifelong_mapf} is a classic planner that does
not learn. It runs in the same loop as the learned controllers, but it
chooses its own tasks and follows paths it plans ahead.
\citet{ma2017lifelong_mapf} give two solvers (Section~\ref{sec:rw:mapf}).
Token Passing is used rather than Token Passing with Task Swaps, because with
task swaps an agent can lose its task to another, while a learned agent keeps
its task once it has one. Token Passing also needs less computation.}
\note{B1, a proposal, see the question.}

\old{\textbf{Token Passing} \citep{ma2017lifelong_mapf} is the only router
that does not learn. The agents share one token that holds every agent's
planned path.} \new{The agents share one token that holds every planned path,
as in Section~\ref{sec:rw:mapf}.} An agent whose path has ended asks for the
token, and the agents ask one at a time in order of index\new{, an order
\citet{ma2017lifelong_mapf} leave open}. The agent takes the waiting task
closest to it in number of moves, among the tasks whose pickup \old{and end
vertices are}\new{$s_j$ and delivery $g_j$ are} not where another agent's path
ends. \old{It then plans a path to the pickup and on to the end vertex by
space-time A$^\star$, around the paths already in the token. Every other agent
is taken to stand where its path ends from then on.} An agent that finds no
such task stays where it is\old{, or moves to the nearest free endpoint if it
stands on the end vertex of a waiting task}\new{. If a waiting task delivers
where it stands, it moves instead to the nearest free endpoint that is no
shelf or delivery point}. \new{Token Passing chooses its own tasks in this way,
so it does not use the assignment rule.} \note{D1, the token is the
Background's, only what is added here stays. B3, the order is ours. C1, the
PF's names. D2, the first of two paragraphs, how a task is chosen.}

\new{The agent then plans the shortest path, in steps, to $s_j$ and on to
$g_j$, keeping clear of the paths already in the token. Every other agent is
taken to stand where its path ends from then on.} Unlike in
\citet{ma2017lifelong_mapf}, where an agent may follow another into the vertex
it is leaving, a path keeps a vertex from one step before an agent arrives
until one step after it leaves. So conflict resolution never has to change a
plan. Two agents therefore never ask for each other's vertex, so Token Passing
never meets the lock at a workstation of Figure~\ref{fig:workstationlock}.
\old{Token Passing chooses its own tasks, so it does not use the assignment
rule.} \old{Its guarantee holds only on well-formed warehouses}\new{It
promises to finish every task when the warehouse is well-formed}, and every
warehouse is built to be one (Appendix~\ref{app:wellformed}). \old{It plays the
role that A$^\star$, ICBS and ICTS play in \citet{vanKnippenberg2021}.}
\note{D2, the second paragraph, how a path is planned. C3, plain words for
space-time A$^\star$, which the thesis does not cite. C4, well-formed is
enough, not needed. B4, dropped.}

\old{Token Passing lets only one agent head for a workstation at a time.}
\new{Only one Token Passing agent heads for a workstation at a time, because
its path ends there and no other task may then deliver to it.} All lines of
an order go to the same workstation, so they are brought one after another,
and a line can wait longer for its turn than its \old{agent drives}\new{trip
takes. The Problem Formulation keeps this wait as a limit of the scope
(Section~\ref{sec:pf:scope})}. Appendix~\ref{app:load:next} measures \old{this}
\new{the} wait. Token Passing's task rule is kept as published, because the
\old{router}\new{planner} is not what this thesis tests. \old{, so its gain
from demand storage is read with this wait in mind.} The learned controllers
take their tasks by the assignment rule and do not have this wait, but they
can lock at a workstation instead (Section~\ref{sec:method:resolution}).
\note{D4, the mechanism in one sentence, then a point back to the scope item.
C2, ``planner''. C5, no ``drive''. How the gain is read moves to the
Experiments.}

\section*{Part 4. Experiments, ``Factors''}

\ldots Communication is varied only for the decentralised controller.
\old{The reference planner of Section~\ref{sec:method:reference} runs under
\mockup{the conditions of issue \#84} as a point of comparison, not as a fourth
controller.}
\new{Token Passing (Section~\ref{sec:method:reference}) is the classic planner
the learned controllers are measured against, not a fourth controller. It runs
under every storage rule at the main load, so its gain can be set next to
theirs. Because it needs no training, it also sets values the design needs. A
load sweep of Token Passing under fixed storage sets the main load of each
fleet (Appendix~\ref{app:load}), and sweeps of it set the storage period and
the pallet size (Appendix~\ref{app:storagecal}). Its gain from demand storage
is read with its workstation wait in mind.} \note{A4. The role by design.
The mockup is filled from Appendices~E and~G.}

\section*{Question}
\begin{enumerate}[nosep]
  \item B1. Token Passing came from the paper that defines lifelong MAPD.
    Is this reason fine: with task swaps an agent can lose its task, while a
    learned agent keeps its task once it has one, and Token Passing needs
    less computation? Or only ``the simpler of the two solvers of Ma et al.''?
\end{enumerate}

\end{document}
